\subsection{Free comparison with global tensor inputs}

Let \(R=\epsilon^{-1}\), \(P=2\epsilon c\le1\), and use
\[
 s=R^2/\rho^2,\qquad q=R\log(\rho/R)+R\tau/2,
 \qquad \eta=Rz/\rho.
\]
The velocity and tensor measures are
\(H=L^2(s^{-3/2}ds\,dq/c\,d\eta)\) and
\(H_T=L^2(s^{-1/2}ds\,dq/c\,d\eta;\operatorname{Sym}_3)\).
The natural lifts \(b\mapsto\sqrt{s}\,b\) and \(T\mapsto sT\)
both multiply norms by \(\sqrt\pi R\). All vector fields are
axisymmetric without swirl, and tensors have no component with exactly
one azimuthal index. Set
\[
 S=s\partial_s,\quad K=2S+\eta\partial_\eta,\quad
 Q_\epsilon=\partial_q-\epsilon(K+1),\quad E=\operatorname{diag}(1,1,0).
\]
Direct Cartesian differentiation gives
\begin{align*}
 A_\epsilon&=-S+\nu s[(Q_\epsilon^2+\partial_\eta^2)I-\epsilon^2E],
 &A_0&=-S+\nu s\Delta_{q,\eta},\\
 \operatorname{div}(\sqrt{s}\,b)
 &=s[(\partial_q-\epsilon K)b_r+\partial_\eta b_z],\\
 (\operatorname{div}_\epsilon T)_r
 &=Q_\epsilon T_{rr}+\partial_\eta T_{rz}-\epsilon T_{\phi\phi},
 &(\operatorname{div}_\epsilon T)_z
 &=Q_\epsilon T_{zr}+\partial_\eta T_{zz}.
\end{align*}
Thus the full tensor divergence, including the hoop traction, is retained.
The domains of \(A_\epsilon\) and its adjoint are the exact Cartesian
domains under these lifts. No boundary is introduced at either age end.
Let \(P_\epsilon\) be the conjugated Cartesian Leray projection and
\(P_0\) the ordinary planar cylinder projection. Write
\[
 T_\epsilon=(\mu-A_\epsilon)^{-1}P_\epsilon,
 \qquad \mathcal S_\epsilon=T_\epsilon s\operatorname{div}_\epsilon.
\]
The tensor map means completion from smooth tensors. The free Cartesian
energy estimate, and its flat age-energy counterpart, give uniformly
for \(0\le\epsilon\le1\) and \(\mu\) in a compact subset of
\(\operatorname{Re}\mu>-1/4\),
\begin{equation}\label{near:free-energy}
 \|b\|_H+\sqrt\nu\|(Q_\epsilon b,\partial_\eta b,
                    \epsilon b_r)\|_{H_T}
 \le C\nu^{-1/2}\|M\|_{H_T},\qquad b=\mathcal S_\epsilon M.
\end{equation}
The flat equation is
\((\partial_s+\mu/s-\nu\Delta)b=P_0\operatorname{div}_0M\).
Pairing it with \(s^{-1/2}\bar b\) gives the coercive constant
\(\operatorname{Re}\mu+1/4\) and proves its completion independently
of comparison. Bounded measurable partitions of \(M\) are therefore
allowed before divergence; their distributional interface terms remain.

\begin{proposition}\label{near:compact-free-comparison}
Let \(\chi\) be supported in \(\sigma\le s\le H_1\),
\(|\eta|\le L_1\), and let \(\chi_1\) equal one on a larger
neighborhood. There is a fixed integer \(k\) such that
\begin{equation}\label{near:compact-free-rates}
 \|\chi(T_\epsilon-T_0)\|_{H\to H}
       \le C\Gamma^k\epsilon^{1/14},\qquad
 \|\chi(\mathcal S_\epsilon-\mathcal S_0)\|_{H_T\to H}
       \le C\Gamma^k\epsilon^{1/50},
\end{equation}
where \(\gamma_*=\inf(\operatorname{Re}\mu+1/4)>0\) over the fixed
parameter set and
\[
 \Gamma=2+\nu+\nu^{-1}+c+c^{-1}+\sigma^{-1}+H_1+L_1
 +\sup|\mu|+\gamma_*^{-1}
 +\|\chi\|_{C^8}+\|\chi_1\|_{C^8}.
\]
The integer is fixed before any window is allowed to depend on
\(\epsilon\). Inputs are global and arbitrary in their indicated spaces.
\end{proposition}

\begin{proof}
First compare the projections on smooth tests. On each nonzero
\(q\)-frequency, a flat divergence-free field and gradient have forms
\((-\psi_\eta,\psi_q)\), \((\varphi_q,\varphi_\eta)\).
Their exact counterparts are
\[
 (-\psi_\eta,(\partial_q-\epsilon K)\psi+\epsilon\psi),
 \qquad ((\partial_q-\epsilon K)\varphi,\varphi_\eta).
\]
The first has zero exact divergence; the second lifts to a Cartesian
gradient. The inverse \(q\)-derivative costs at most \(c/(2\pi)\).
Also, for the flat Fourier multiplier,
\(|\partial_\xi^jP_0(k,\xi)|\le C_j|k|^{-j}\), \(j=1,2\),
so one or two transverse moments of the projected field cost only
fixed powers of \(c\). All projection operations commute with \(S\).

For the zero frequency, the radial component is divergence-free and the
axial component is a gradient. Take the corresponding primitives to
vanish at \(\eta=0\), and cut them off at \(|\eta|\sim L_0\).
Their suprema, including one \(S\) derivative, are bounded by
\(C\|\langle\eta\rangle f\|_2\) and
\(C\|\langle\eta\rangle Sf\|_2\).
Cutoff differentiation costs \(L_0^{-1/2}\), geometric correction
costs \(\epsilon L_0^{1/2}\), and the omitted field tail costs
\(L_0^{-1}\). Choose \(L_0=\epsilon^{-1}\).
Distances to the exact gradient and divergence-free subspaces, both closed
orthogonal subspaces of \(H\), now give
\begin{equation}\label{near:leray-test}
 \|(P_\epsilon-P_0)f\|_H
 \le p(1+c+c^{-1})\sqrt\epsilon
              \sum_{a,j=0}^1\|\eta^jS^af\|_H.
\end{equation}
Initially use finite Fourier sums, compact log-age support and Schwartz
transverse profiles. The exact corrections lift to annular fields.
Cutoff, mollification and Fourier truncation are dense in the displayed
norm, after \(x=\log s\), \(u=e^{-x/4}f(e^x)\).
Thus \eqref{near:leray-test} holds on its full stated domain and retains
the zero mode.

Next obtain uniform local regularity of the outputs.
The unitary lift conjugates \(\mathcal T-3/4\) to \(S-1/4\).
The Cartesian estimate \eqref{free:half-drift} therefore gives
\[
 \||S-1/4|^{1/2}\mathcal S_\epsilon M\|_H
           \le C\nu^{-1/2}\|M\|_{H_T}\qquad(\epsilon>0).
\]
At \(\epsilon=0\), a compact-age cutoff \(u\) of the flat solution
satisfies \(u\in L_s^2H^1\), \(u_s\in L_s^2H^{-1}\) by
\eqref{near:free-energy} and its equation. Fourier Cauchy--Schwarz gives
\(\||\partial_s|^{1/2}u\|_2^2\le
\|u_s\|_{L_s^2H^{-1}}\|u\|_{L_s^2H^1}\).
All local constants are polynomial in \(\Gamma\).

In log age the conjugate of \(Q_\epsilon\), shifted by \(\epsilon\),
is the skew operator
\(\widetilde Q_\epsilon=\partial_q-2\epsilon\partial_x
                     -\epsilon(\eta\partial_\eta+1/2)\).
Its three constituent flows commute. Multiplication by a smooth cutoff
preserves the half age derivative: its difference quotient splits into
the cutoff times the quotient of the function and its first derivative
times a translated function. Both terms have the stated polynomial bound.
For a compactly supported
cutoff \(f\) of \(e^{-x/4}b(e^x)\), their unitary flow identities give
\[
 \|f(q+h)-f(q)\|_2
 \le |h|\|\widetilde Q_\epsilon f\|_2
 +C|\epsilon h|^{1/2}\||\partial_x|^{1/2}f\|_2
 +|\epsilon h|\|(\eta\partial_\eta+1/2)f\|_2.
\]
For \(|h|\le1\), square and integrate against \(|h|^{-3/2}\);
for larger \(h\), use \(2\|f\|_2\). The half age derivative
and full transverse derivative already have bounds. Directional Fourier
inequalities and the smooth log-age change prove
\begin{equation}\label{near:local-smoothing}
 \|\chi\mathcal S_\epsilon M\|_{H^{1/4}(s,q,\eta)}
       \le C\Gamma^k\|M\|_{H_T}.
\end{equation}
For a body source the graph controls \(Sb\), and the gradient controls
\(Q_\epsilon b,\partial_\eta b\). Since
\(\partial_s=s^{-1}S\) and
\(\partial_q=Q_\epsilon+\epsilon(2S+\eta\partial_\eta+1)\),
the body output has a full local \(H^1\) bound.
The flat body graph follows directly by cutting off its heat equation
on dyadic intervals \([a/2,4a]\): the Fourier heat estimate controls
\(\nu s\Delta b\), the equation controls \(Sb\), and finite overlap
sums these bounds. Periodic extension over many cells followed by division
by their number gives constants independent of the number of cells.

It remains to compare smooth adjoint observations globally. For a smooth
test \(z\) in the observation box, put \(h=P_0z\) and \(v=T_0^*z\).
Exactly,
\[
 v(s)=\int_s^\infty(s/t)^{\bar\mu+1/2}
                e^{\nu(t-s)\Delta}h(t)\,\frac{dt}{t}.
\]
It vanishes for \(s>H_1\). With \(a=\bar\mu+1/2\),
\[
 Sv=av-\nu s\Delta v-h,\qquad
 S^2v=a^2v-\nu s(2a+1)\Delta v+\nu^2s^2\Delta^2v
                         -ah+\nu s\Delta h-Sh.
\]
In log age the integral is a heat contraction followed by translation,
with scalar kernel \(e^{-(\bar\mu+1/4)r}\), \(r>0\).
Minkowski bounds it by the inverse coercivity gap. The identities
\[
 \eta e^{t\Delta}=e^{t\Delta}(\eta-2t\partial_\eta),\qquad
 \eta^2e^{t\Delta}=e^{t\Delta}
       (\eta^2-4t\eta\partial_\eta-2t+4t^2\partial_\eta^2)
\]
control two transverse moments since the heat time is at most \(\nu H_1\).
The multiplier estimate used in \eqref{near:leray-test} transfers these
moments through \(P_0\). The preceding two age identities therefore give,
with \(\Lambda=K+1\),
\begin{equation}\label{near:adjoint-polynomial}
 \|\nu s\Lambda\partial_qv\|_H+\|\nu s\Lambda^2v\|_H+\|\nu sEv\|_H
 +\|\eta\partial_\eta v\|_H
 +\sum_{a,j=0}^1\|\eta^jS^av\|_H
 \le C\Gamma^k\|z\|_{H^6}.
\end{equation}
Indeed \(\Lambda^2=4S^2+4S\eta\partial_\eta+4S+
\eta^2\partial_\eta^2+3\eta\partial_\eta+1\).
At most two age differentiations, two moment commutations and six test
derivatives occur. Below the test support all norms are integrable by
\(\int_0^\sigma s^{2\operatorname{Re}\mu-1/2}ds<\infty\).
Thus the degree \(k\) is fixed independently of later windows.

The exact generator difference is
\(A_\epsilon-A_0=\nu s[-2\epsilon\Lambda\partial_q+
\epsilon^2\Lambda^2-\epsilon^2E]\), also for adjoints. If
\(w=T_\epsilon^*z-v\), then
\[
 (\bar\mu-A_\epsilon^*)w
 =(P_\epsilon-P_0)z+(A_\epsilon^*-A_0^*)v.
\]
The right side has \(H\) norm at most
\(C\Gamma^k\sqrt\epsilon\|z\|_{H^6}\).
The exact componentwise free adjoint energy estimate bounds \(w\)
and its physical gradient by this quantity (with a polynomial viscosity
cost). There is no projection before the second source term.
The domain application is legitimate because
\eqref{near:adjoint-polynomial} puts \(v\) in the actual maximal adjoint
graph; its lift vanishes near the axis, without any assumption that
\(v\) is divergence-free for finite \(\epsilon\).

Let \(G_\epsilon\) be the negative weighted symmetric gradient adjoint
to \(s\operatorname{div}_\epsilon\). Then
\[
 (\mathcal S_\epsilon^*-\mathcal S_0^*)z
   =G_\epsilon w+(G_\epsilon-G_0)v.
\]
The divergence formula at the start shows that the second term is
\(\epsilon\) times \(Sv,\eta\partial_\eta v,v\) and frame terms.
Consequently
\begin{equation}\label{near:smooth-adjoint-comparison}
 \|(T_\epsilon^*-T_0^*)z\|_H
 +\|(\mathcal S_\epsilon^*-\mathcal S_0^*)z\|_{H_T}
 \le C\Gamma^k\sqrt\epsilon\|z\|_{H^6}.
\end{equation}

Place the output support in a larger periodic box in \(s,\eta\),
retaining the original normalized \(q\) measure. Fourier projection
to frequencies at most \(N\), and the local \(H^1\) and
\(H^{1/4}\) bounds, give high-frequency errors \(C\Gamma^kN^{-1}\)
and \(C\Gamma^kN^{-1/4}\). A unit low-frequency dual vector corresponds
to the compact smooth \(H\)-test \(s^{3/2}\chi F_Nh\), with
\(H^6\) norm at most \(C\Gamma^kN^6\). Thus the full operator
differences are at most
\[
 C\Gamma^k(N^{-1}+\sqrt\epsilon N^6),\qquad
 C\Gamma^k(N^{-1/4}+\sqrt\epsilon N^6).
\]
Choose \(N=\epsilon^{-1/14}\) and \(N=\epsilon^{-2/25}\), respectively.
No Fourier mode-count factor is used. This proves
\eqref{near:compact-free-rates}.
\end{proof}

The argument uses only normalized phase averages. The imported Cartesian
graph bound is obtained by averaging fixed-size heat windows over their
centers, so it has no inverse-period loss. Every remaining operation
above is a spatial or age identity. Hence the comparison remains valid
when \(P=2\epsilon\nu^{-2}\) tends to zero. Exact Cartesian divergence,
the pressure projection, the zero cylinder mode and tensor interface
terms have all been retained.
